%% Evaluation design and E1--E3. Section labels follow the historical source.
\section{Evaluation}
\label{sec:eval}
The evaluation follows a document from ingestion and organisation to
retrieval and use. E1, E2 and E4 test the managed state and its lifecycle
(RQ1). E3 tests whether catalogue routing selects the intended collection
(RQ2), and E5 measures what each retrieval configuration puts into its
evidence pack (RQ2--RQ3). E8 tests how verifier-based ordering changes the
top five (RQ4). E6 and E7 examine answer traces, refusals and execution
costs (RQ5), and E9 reports operational checks across the workflow.
Tables~S9--S10 and S17 record the protocols and provenance.

\subsection{Experimental design and measures}
\label{sec:eval-setup}

\textbf{Corpora and questions.} The industrial corpus contains 42 documents
in five libraries, including twenty near-duplicate telemetry reports. It
has 16 questions with annotated document fragments. Q16 shares a token with
its gold filename. A research-paper campaign sampled 100 papers from arXiv
across about 25 domains; ten questions have gold papers annotated by their
source identifiers, which are absent from the saved question texts.
A November
metadata recheck enumerated 328 parts in five libraries, including six parts
from non-target documents. The
cross-lingual campaign used 400 Wikipedia pages and 40 extractive queries per
language (English, Chinese and Japanese). Its ingestion script placed pages
in predefined language libraries and checked the resulting listings.
The routing and content benchmark used 138 documents stored as 879 parts in
ten domains and 72 in-domain questions. The released traces for complete
arms also contain four out-of-domain probes (R1--R4), reported separately.
The benchmark excludes the two cross-domain placeholders X1 and X2. The
follow-up used the same 72 questions after the corpus grew to 906 parts.
The questions and gold annotations were developer-authored. A post-hoc
check lowercased each question and its nonempty annotated target string,
then tested whether the entire target was a substring of the question,
without other normalization. This identified overlap in 29 of 72 questions;
the remaining 43 did not match under this test. The check describes the
questions and does not rescore the returned packs.

\textbf{Measures and controls.} Document-level Hit@5, Recall@5 and
MRR@5 are evaluated on the five returned documents, with reciprocal rank set
to zero when the top five contain no gold item \citep{craswell-msmarco-2021}.
The route check tests whether the script selected the labelled collection.
Target-string containment records whether the annotated string occurs
anywhere in the returned evidence pack, including in a sibling document.
The answer grids check for a gold paper identifier in the answer or trace,
for at least $\lceil 3n/5\rceil$ of $n$ annotated terms (strict matches or,
in the alias-aware check, accepted variants), and for a document-style
citation marker. These lexical checks measure the presence of expected
text and citation formatting. Semantic assessment and graded relevance
would require distinct criteria, as in RAG evaluation frameworks
\citep{es-ragas-2024,saad-falcon-ares-2024} and cumulated-gain measures
\citep{jarvelin-cumulated-gain-2002}.

Wilson intervals describe the observed question sets. Paired exact
McNemar tests are two-sided and exploratory, with unadjusted $p$ values.
Each agent arm used one model and one run per condition, with transport
retries treated as part of that run. The answer grids used the platform chat
API with the Claude Code engine. A separate four-question case compared
Claude Code with oh-my-pi using \texttt{ustc/deepseek-flash} in November
2026. Two corpus-integrity defects were repaired before measurement.

\subsection{E1: Acquisition and parsing}
\label{sec:e1}
The campaign targeted 100 papers. In the November 14--15 disk inventory,
every target had at least one stored part; the five libraries contained 328
parts in total, including six from non-target documents. All 328 had
metadata descriptions of at least 60 characters without the recorded
placeholder pattern. This establishes neither complete paper coverage nor
retrieval quality.

Earlier records report 91 names in an ingest-state list and 79 documents
enumerated and read from a derived library. These are not matched target
sets or stages of a 100-paper funnel, and the records do not establish a
reliable date or individual failure stages. An initial November probe also
recorded misses. After forced reindexing, explicit indexing and a probe-tool
correction, a filename-stem short-query diagnostic found all 100 target
identifiers in returned names. However, 72 searches were unrestricted by
library and 45 returned a chunk-copy name first; this is not a matched
before--after retrieval comparison. A separate vector check searched one
preassigned library for each of ten questions and ranked the gold document
first in all ten. The judge-input loader could fall back to the known
source identifier; the retained summary does not say when it did so.
This check does not measure end-to-end routing, judging or evidence packing.

The Wikipedia ingestion check found 350 of 400 expected titles in the
preassigned English library and 400 of 400 in each of the Chinese and
Japanese libraries. Its separate self-retrieval check used the first
three index results. These tests exercise document placement and
index access for known inputs. In a 50-paper campaign spanning 26
domains, all 50 papers passed the parsing and verification checks.
The run recorded 41 negative tag flags without blocking ingestion.

\subsection{E2: Organisation, consistency, and management operations}
\label{sec:e2}
A rename in a test library exposed stale index state. The index returned
18 chunks where nine were expected, producing duplicate results. Forced
reindexing restored the expected nine. A separate five-library inventory
recorded 322 parts; its archived content snippets were nonempty for all
322, but retained at most the first 1,000 characters per part. A 33-part
sample passed checks across files, tree entries and
library metadata. A three-level
test library supported parent-to-child moves, listing, reading and
judging. Of the 16 lifecycle scenarios, the create, move, delete and
forced reindex checks passed, and the task-status and update checks
failed. The campaign left merge and split outside the controlled tests.

Near-duplicate review flagged 35 candidate groups at a threshold of~0.90
(scores 0.90--0.975). Review found no true duplicates and removed no
documents (Table~S9). The tag check passed 165 parts and exposed orphan
tags and one cross-library registry leak. A document move in the 658-node
graph added four tag edges (4,675 to 4,679), recording the change in the
navigation view.

\subsection{E3: Paired route-card audit: the description-layer contrast}
\label{sec:routing-audit}
In the frozen audit, the script selected the labelled collection for all
72 in-domain questions using description-bearing route cards (A0), and for
37 using nominally name-only cards (A3). The 35 changed decisions all
favoured A0 (Figure~\ref{fig:routing}a, exact McNemar $p=5.82\times10^{-11}$).
In both later records, A0 selected the labelled collection in 72 cases and
A3 in 49. These are contrasts between route-card configurations. A3 could fall back
to descriptions without logging the fallback, some questions share
titles with the descriptions, and A3 can search outside the evaluation
scope. Among the 72 questions, 52 were expanded automatically, and
47 came from three domains.

In the frozen record, the returned packs contained the target string in
37 A0 and 33 A3 cases (10 A0-only, six A3-only, $p=0.4545$). The mechanical full
scan located all 72 targets (Figure~\ref{fig:routing}b), including the
35 missed by A0 ($p=5.82\times10^{-11}$). It issued 64{,}584 tool calls
across these questions, averaging 897 calls and 24.3 MB of recorded text
per question. The containment--text-volume plot in
Figure~\ref{fig:routing}c shows the cost of this reference point.
Section~\ref{sec:ladder-opt} compares the evidence packs across
retrieval configurations.

\begin{figure}[!t]
\centering
\includegraphics[width=\columnwidth]{fig_opt_route}
\caption{Frozen route-card and evidence-pack diagnostics on 72 in-domain
questions. (a) Labelled-collection selection with Wilson 95\% intervals.
(b) Annotated target-string containment in returned text. The read-all
arm performs a mechanical substring scan and generates no answers.
(c) Frozen trace-recorded text per question (MB, log scale) and
containment. The A0f and PD markers are offset horizontally by 8~pt
for legibility. The source data retain their original coordinates.
Recorded text and tool time measure the retrieval work in the trace.}
\label{fig:routing}
\end{figure}
